\section{ViT：把图像变成序列意味着什么}

有了数据规模与 transfer protocol，讲座才进入 ViT。其研究假设不是“attention 一定优于 convolution”，而是：当数据和 compute 足够时，能否减少 locality 与 translation equivariance 等硬编码视觉先验，让模型从数据中学出空间结构？因此 ViT 的简单性既是优点，也是对规模的要求。

\subsection{Let the data speak：减少手工视觉结构}

在 CNN 中，local receptive field、weight sharing 与层级下采样形成强归纳偏置。ViT 保留更少的二维结构：先把图像分块并线性投影，再用通用 Transformer encoder 处理 token。研究问题从“如何设计更好的卷积模块”变成“数据能否补回被移除的先验”。

\lecturefigure{slide-16-letting-data-speak.jpg}{“ViT: letting the data speak”把架构简化视为规模化实验。}{官方视频 25:12。}

\begin{knowledgebox}{读图：更少先验不是更少假设}
Patch size、position embedding、class token、augmentation、optimizer 与 dataset composition 仍是强设计选择。所谓 letting the data speak，是相对减少 convolutional locality，而不是消除 inductive bias。更弱先验通常提高表示灵活性，也扩大 sample complexity 和 optimization burden。
\end{knowledgebox}

\lecturefigure{slide-17-cnn-vs-transformer.jpg}{CNN 通过卷积层级编码视觉结构；Transformer 在语言中已证明可扩展，ViT 询问它能否直接用于图像。}{官方视频 26:24。}

\begin{knowledgebox}{读图：比较的是信息路径与参数共享方式}
卷积让邻近像素先局部交互，并跨层扩大 receptive field；self-attention 允许任意 patch 在一层中建立内容依赖连接。CNN 的局部性更 data-efficient；attention 的全局交互更灵活，但原始复杂度随 token 数二次增长。二者并非“局部 vs 全局”的绝对二分，深 CNN 也有全局 receptive field，ViT 也会学习局部 pattern。
\end{knowledgebox}

\subsection{Patch tokenization 与 class token}

设输入图像 $x\in\mathbb{R}^{H\times W\times C}$，patch 边长为 $P$，则 token 数为 $N=HW/P^2$。每个 patch 展平为 $P^2C$ 维向量，经线性矩阵 $E$ 投影到 $D$ 维。再前置可学习 class token，并加入 position embedding：
\[
z_0=[x_{\mathrm{cls}};x_p^1E;x_p^2E;\ldots;x_p^NE]+E_{\mathrm{pos}}.
\]
这里 $x_p^i$ 是第 $i$ 个 patch，$x_{\mathrm{cls}}$ 用于聚合分类信息，$E_{\mathrm{pos}}$ 提供 token 的空间顺序。

\lecturefigure{slide-18-vit-architecture.jpg}{ViT 将 patch embedding、position embedding 与 class token 输入标准 Transformer encoder。}{官方视频 28:32；Dosovitskiy et al. 2020。}

\begin{knowledgebox}{读图：架构图从下往上读}
最下方图像被切成固定 patch；线性投影把每个 patch 变成 token；position embedding 防止 permutation invariance 丢失空间；encoder 交替执行 multi-head self-attention 与 MLP；最终 class token 进入分类 head。图中“标准 Transformer”仍包含 LayerNorm、residual connection 与多层堆叠，并非只有 attention 一步。
\end{knowledgebox}

\begin{importantbox}{Self-attention 的计算代价}
对 $N$ 个 token、hidden size $D$，attention score $QK^{\top}$ 约需 $O(N^2D)$ 计算和 $O(N^2)$ 中间存储。因为 $N=HW/P^2$，把 patch 边长减半会让 token 数变为四倍，attention matrix 面积变为十六倍。高分辨率视觉任务因此不能只看参数量，还要看 token 数与 activation memory。
\end{importantbox}

\begin{lstlisting}[caption={Patchify and prepend a class token}]
def vit_tokens(image, patch_size, projection, cls_token, pos_embed):
    patches = split_into_patches(image, patch_size)
    tokens = flatten(patches) @ projection
    tokens = concatenate([cls_token, tokens], axis=0)
    return tokens + pos_embed
\end{lstlisting}

\subsection{本章小结}

ViT 的核心实验是弱化卷积先验，而不是证明图像“本质上就是语言”。Patchify 决定 token 数与 compute，position embedding 恢复顺序，class token 提供聚合接口；简单架构把更多负担交给数据、优化和规模。

